\documentclass[10pt,a4paper]{article}

% --- PACKAGES ---
\usepackage[utf8]{inputenc}
\usepackage[T1]{fontenc}
\usepackage[scaled=0.92]{helvet}
\usepackage{mathpazo}
\usepackage[margin=1cm, top=1.2cm, bottom=1.2cm]{geometry}
\usepackage{amsmath,amssymb}
\usepackage{booktabs}
\usepackage{tabularx}
\usepackage{multicol}
\usepackage{enumitem}
\usepackage{titlesec}
\usepackage{fancyhdr}
\usepackage{xcolor}
\usepackage{listings}
\usepackage[hidelinks]{hyperref}
\usepackage{setspace}

% --- SETUP ---
\setlength{\parindent}{0pt}
\setlength{\parskip}{0.3em}
\setlist[itemize]{leftmargin=1em, noitemsep, topsep=1pt, parsep=0pt, itemsep=0pt}
\setlist[enumerate]{leftmargin=1.2em, noitemsep, topsep=1pt, parsep=0pt, itemsep=0pt}
\setstretch{1.02}

% --- COLORS & STYLES ---
\definecolor{primary}{HTML}{1F3A5F}
\definecolor{codebg}{HTML}{F4F5F7}
\definecolor{warnfg}{HTML}{9C4221}

\titleformat{\section}{\normalfont\large\bfseries\sffamily\color{primary}}{}{0em}{}[\vspace{-0.5em}\rule{\linewidth}{0.5pt}]
\titlespacing*{\section}{0pt}{1.2ex}{0.8ex}

\lstset{
    basicstyle=\ttfamily\footnotesize,
    keywordstyle=\color{primary}\bfseries,
    backgroundcolor=\color{codebg},
    frame=none,
    breaklines=true,
    aboveskip=3pt,
    belowskip=3pt,
    xleftmargin=0.4em,
    xrightmargin=0.4em,
    columns=fullflexible,
    keepspaces=true,
    showstringspaces=false
}

\pagestyle{fancy}
\fancyhf{}
\fancyhead[L]{\small\sffamily\bfseries\color{primary} CHEAT SHEET BASIS DATA LENGKAP}
\fancyhead[R]{\small\sffamily Muhammad Riduwan Khafidi}
\renewcommand{\headrulewidth}{0.4pt}
\renewcommand{\headrule}{\hbox to\headwidth{\color{gray!50}\leaders\hrule height \headrulewidth\hfill}}
\setlength{\headheight}{22pt}

\begin{document}

\vspace{-1em}
\begin{multicols*}{2}

	% ==========================================
	\section*{1. Konsep Dasar \& ERD}

	\textbf{Kamus:} Data = fakta mentah. Database = kumpulan data terorganisir. DBMS = software manajemen (MySQL, MariaDB, Oracle, PostgreSQL). RDBMS = DBMS berbasis relasi tabel.

	\textbf{Notasi ERD (Chen):}
	\begin{tabularx}{\linewidth}{@{}l X@{}}
		\textbf{Kotak}         & Entitas (objek nyata)     \\
		\textbf{Elips}         & Atribut (properti)        \\
		\textbf{Elips Ganda}   & Atribut Multivalued       \\
		\textbf{Elips Putus}   & Atribut Turunan (derived) \\
		\textbf{Belah Ketupat} & Relasi                    \\
		\textbf{Kotak Ganda}   & Entitas Lemah             \\
		\textbf{Garis Bawah}   & Primary Key               \\
		\textbf{Garis Ganda}   & Partisipasi Total         \\
	\end{tabularx}

	\textbf{Kardinalitas:}
	\begin{itemize}
		\item \textbf{1:1} — Dept punya 1 Manajer
		\item \textbf{1:N} — Dept punya banyak Pegawai
		\item \textbf{N:M} — Mhs ambil banyak Matkul
	\end{itemize}

	\textbf{EERD (Lanjutan):}
	\begin{itemize}
		\item \textbf{Generalization:} Segitiga IS-A. Pegawai → Dosen, Staff.
		\item \textbf{Aggregation:} Relasi jadi entitas baru.
	\end{itemize}

	% ==========================================
	\section*{2. Normalisasi}

	\textbf{Tujuan:} hilangkan redundansi, cegah anomali insert/update/delete.

	\begin{itemize}
		\item \textbf{1NF} — Setiap sel \textit{atomik}, tidak boleh array/multi-nilai.
		\item \textbf{2NF} — 1NF + tidak ada \textit{partial dependency} (non-PK tidak bergantung sebagian PK komposit).
		\item \textbf{3NF} — 2NF + tidak ada \textit{transitive dependency} (non-PK tidak bergantung non-PK lain).
		\item \textbf{BCNF} — 3NF + setiap determinant adalah candidate key.
		\item \textbf{4NF} — BCNF + tidak ada \textit{multi-valued dependency}.
	\end{itemize}

	\textbf{Mnemonic Codd:} \textit{``The Key, The Whole Key, and Nothing But The Key.''}

	% ==========================================
	\section*{3. Konversi ERD $\to$ Tabel}

	\begin{enumerate}
		\item \textbf{Entitas Kuat} $\to$ Tabel. PK entitas $\to$ PK tabel.
		\item \textbf{Entitas Lemah} $\to$ Tabel. PK entitas kuat masuk sebagai FK + PK komposit. Beri \texttt{ON DELETE CASCADE}.
		\item \textbf{Relasi 1:1} $\to$ Pilih salah satu tabel, taruh PK lainnya sebagai FK.
		\item \textbf{Relasi 1:N} $\to$ PK sisi ``1'' masuk sebagai FK di sisi ``N''.
		\item \textbf{Relasi N:M} $\to$ \textbf{WAJIB} tabel baru (junction). PK komposit dari kedua FK.
		\item \textbf{Atribut Multivalued} $\to$ Tabel baru.
		\item \textbf{Atribut Turunan} $\to$ JANGAN disimpan, hitung saat query.
	\end{enumerate}

	% ==========================================
	\section*{4. DDL — Struktur}

	\textbf{Database:}
	\begin{lstlisting}[language=SQL]
CREATE DATABASE [IF NOT EXISTS] db;
USE db;
SHOW DATABASES;
DROP DATABASE [IF EXISTS] db;
\end{lstlisting}

	\textbf{Tabel:}
	\begin{lstlisting}[language=SQL]
CREATE TABLE [IF NOT EXISTS] tbl (
  kolom1 TIPE [constraints],
  kolom2 TIPE [constraints],
  PRIMARY KEY (kolom),
  FOREIGN KEY (kolom)
    REFERENCES tbl_lain(kolom)
    ON DELETE CASCADE
);
\end{lstlisting}

	\textbf{Alter:}
	\begin{lstlisting}[language=SQL]
ALTER TABLE t ADD COLUMN c TIPE;
ALTER TABLE t DROP COLUMN c;
ALTER TABLE t MODIFY COLUMN c TIPE;
ALTER TABLE t CHANGE c1 c2 TIPE;
ALTER TABLE t RENAME TO t_baru;
ALTER TABLE t ADD PRIMARY KEY (c);
ALTER TABLE t DROP PRIMARY KEY;
ALTER TABLE t ADD INDEX idx (c);
\end{lstlisting}

	\textbf{Hapus:}
	\begin{lstlisting}[language=SQL]
TRUNCATE TABLE t;  -- hapus baris saja
DROP TABLE [IF EXISTS] t;  -- hapus semua
\end{lstlisting}

	% ==========================================
	\section*{5. Tipe Data}

	\textbf{Integer:} \texttt{TINYINT} (1B), \texttt{SMALLINT} (2B), \texttt{MEDIUMINT} (3B), \texttt{INT} (4B), \texttt{BIGINT} (8B). Tambah \texttt{UNSIGNED} untuk non-negatif.

	\textbf{Desimal:}
	\begin{itemize}
		\item \texttt{DECIMAL(p,s)} — EXACT, untuk UANG
		\item \texttt{FLOAT} — approx 4B, ilmiah
		\item \texttt{DOUBLE} — approx 8B, ilmiah
	\end{itemize}

	\textbf{String:}
	\begin{itemize}
		\item \texttt{CHAR(n)} — fixed, selalu n byte. Kode negara, kode pos.
		\item \texttt{VARCHAR(n)} — variabel. Nama, email, alamat.
		\item \texttt{TEXT} — panjang. Artikel, deskripsi.
		\item \texttt{ENUM('a','b')} — pilihan terbatas.
	\end{itemize}

	\textbf{Waktu:}
	\begin{itemize}
		\item \texttt{DATE} — YYYY-MM-DD
		\item \texttt{TIME} — HH:MM:SS
		\item \texttt{DATETIME} — gabungan
		\item \texttt{TIMESTAMP} — UTC, bisa auto-update
		\item \texttt{YEAR} — 4 digit
	\end{itemize}

	\textbf{Lain:} \texttt{BOOLEAN} (alias \texttt{TINYINT(1)}), \texttt{BINARY(n)}, \texttt{VARBINARY(n)}, \texttt{BLOB}, \texttt{JSON}.

	% ==========================================
	\section*{6. Constraints}

	\begin{tabularx}{\linewidth}{@{}l X@{}}
		\texttt{NOT NULL}        & Wajib diisi                      \\
		\texttt{UNIQUE}          & Tidak boleh duplikat, boleh NULL \\
		\texttt{PRIMARY KEY}     & Unik + tidak null, 1 per tabel   \\
		\texttt{FOREIGN KEY}     & Referensi ke PK tabel lain       \\
		\texttt{CHECK}           & Nilai harus penuhi kondisi       \\
		\texttt{DEFAULT}         & Nilai bawaan                     \\
		\texttt{AUTO\_INCREMENT} & Isi otomatis naik                \\
	\end{tabularx}

	% ==========================================
	\section*{7. Index}

	\begin{lstlisting}[language=SQL]
CREATE INDEX idx_nama ON tbl (kolom);
CREATE UNIQUE INDEX idx ON tbl (kolom);
DROP INDEX idx ON tbl;
\end{lstlisting}

	\textbf{Buat index jika:} kolom sering di \texttt{WHERE}, \texttt{JOIN}, \texttt{ORDER BY}, selektivitas tinggi.

	\textbf{JANGAN buat jika:} tabel kecil, kolom sering di-\texttt{UPDATE}, kolom dengan sedikit nilai unik.

	% ==========================================
	\section*{8. DML — SELECT \& WHERE}

	\begin{lstlisting}[language=SQL]
SELECT kolom FROM tbl;
SELECT DISTINCT kolom FROM tbl;
SELECT * FROM tbl LIMIT 10 OFFSET 20;
SELECT kolom AS alias FROM tbl;
SELECT * FROM tbl ORDER BY c1 ASC, c2 DESC;
\end{lstlisting}

	\textbf{Operator WHERE:}
	\begin{itemize}
		\item Logika: \texttt{AND}, \texttt{OR}, \texttt{NOT}
		\item Banding: \texttt{=, >, >=, <, <=, <>}
		\item \texttt{LIKE}: \texttt{\%} (0+ huruf), \texttt{\_} (1 huruf)
		\item \texttt{BETWEEN a AND b}: inklusif
		\item \texttt{IN (a,b,c)}: salah satu
		\item \texttt{IS NULL} / \texttt{IS NOT NULL}
		\item \texttt{EXISTS (subquery)}
	\end{itemize}

	\textbf{Contoh LIKE:} \texttt{'Fr\%'} (awal Fr), \texttt{'\%re'} (akhir re), \texttt{'\_re'} (2 huruf + re).

	% ==========================================
	\section*{9. Fungsi Built-in}

	\textbf{Aggregate:} \texttt{COUNT}, \texttt{SUM}, \texttt{AVG}, \texttt{MIN}, \texttt{MAX}, \texttt{GROUP\_CONCAT}.

	\textbf{String:}
	\begin{lstlisting}[language=SQL]
CONCAT(a,b)  UPPER(s)  LOWER(s)
SUBSTRING(s,1,3)  LENGTH(s)
TRIM(s)  REPLACE(s,'x','y')
LOCATE('x', s)
\end{lstlisting}

	\textbf{Date:}
	\begin{lstlisting}[language=SQL]
NOW()  CURDATE()  CURTIME()
YEAR(d)  MONTH(d)  DAY(d)
DATEDIFF(d1,d2)  -- selisih hari
DATE_ADD(d, INTERVAL 7 DAY)
DATE_FORMAT(d, '%Y-%m-%d')
\end{lstlisting}

	\textbf{Numeric:}
	\begin{lstlisting}[language=SQL]
ROUND(x,n)  CEIL(x)  FLOOR(x)
ABS(x)  MOD(a,b)  POWER(x,2)
\end{lstlisting}

	% ==========================================
	\section*{10. Aggregate \& GROUP BY}

	\begin{lstlisting}[language=SQL]
SELECT branch, COUNT(*), AVG(balance)
FROM Account
GROUP BY branch
HAVING AVG(balance) > 1000
ORDER BY branch;
\end{lstlisting}

	\textbf{Aturan:} semua kolom di \texttt{SELECT} yang bukan aggregate \textbf{WAJIB} ada di \texttt{GROUP BY}.

	\textbf{WHERE vs HAVING:} \texttt{WHERE} filter sebelum grouping. \texttt{HAVING} filter setelah grouping.

	% ==========================================
	\section*{11. JOIN}

	\textbf{Jenis:}
	\begin{itemize}
		\item \texttt{INNER JOIN} — irisan (cocok di keduanya)
		\item \texttt{LEFT JOIN} — semua kiri + cocok kanan
		\item \texttt{RIGHT JOIN} — semua kanan + cocok kiri
		\item \texttt{FULL OUTER} — semua kedua sisi
		\item \texttt{CROSS JOIN} — perkalian cartesian
		\item \texttt{SELF JOIN} — tabel dengan dirinya
	\end{itemize}

	\begin{lstlisting}[language=SQL]
SELECT A.Owner, D.Amount
FROM Account A
INNER JOIN Deposit D
  ON A.Number = D.Account
WHERE D.Amount > 1000;
\end{lstlisting}

	\textbf{Gaya lama (implicit):}
	\begin{lstlisting}[language=SQL]
FROM Account A, Deposit D
WHERE A.Number = D.Account
\end{lstlisting}

	% ==========================================
	\section*{12. Subquery}

	\textbf{Scalar} (1 nilai):
	\begin{lstlisting}[language=SQL]
SELECT nama FROM customer
WHERE balance > (SELECT AVG(balance)
                 FROM customer);
\end{lstlisting}

	\textbf{IN / NOT IN}:
	\begin{lstlisting}[language=SQL]
SELECT * FROM customer
WHERE cust_name IN
  (SELECT cust_name FROM borrower);
\end{lstlisting}

	\textbf{EXISTS} (korelasi):
	\begin{lstlisting}[language=SQL]
SELECT * FROM customer c
WHERE EXISTS
  (SELECT 1 FROM borrower b
   WHERE b.cust_name = c.cust_name);
\end{lstlisting}

	% ==========================================
	\section*{13. CRUD}

	\begin{lstlisting}[language=SQL]
INSERT INTO t (c1,c2) VALUES (v1,v2);
INSERT INTO t VALUES (v1,v2,v3);
INSERT INTO t SELECT ...;
\end{lstlisting}

	\begin{lstlisting}[language=SQL]
UPDATE t SET c1=v1 WHERE kondisi;
UPDATE t SET balance = CASE
  WHEN balance > 1000
    THEN balance * 1.06
  ELSE balance * 1.05 END;
\end{lstlisting}

	\begin{lstlisting}[language=SQL]
DELETE FROM t WHERE kondisi;
DELETE FROM t;  -- semua baris
\end{lstlisting}

	% ==========================================
	\section*{14. Transaction \& ACID}

	\begin{lstlisting}[language=SQL]
START TRANSACTION;
  ...perintah...
COMMIT;    -- simpan permanen
ROLLBACK;  -- batalkan
SAVEPOINT sp1;
ROLLBACK TO sp1;
\end{lstlisting}

	\textbf{ACID:}
	\begin{itemize}
		\item \textbf{Atomicity} — semua atau tidak
		\item \textbf{Consistency} — tetap valid
		\item \textbf{Isolation} — tidak saling ganggu
		\item \textbf{Durability} — permanen setelah commit
	\end{itemize}

	\textbf{Isolation Level} (longgar $\to$ ketat):
	\begin{enumerate}
		\item \texttt{READ UNCOMMITTED} (dirty read)
		\item \texttt{READ COMMITTED} (non-repeatable)
		\item \texttt{REPEATABLE READ} (phantom) — default MySQL
		\item \texttt{SERIALIZABLE} (paling ketat)
	\end{enumerate}

	% ==========================================
	\section*{15. Views}

	\begin{lstlisting}[language=SQL]
CREATE VIEW v AS SELECT ...;
CREATE OR REPLACE VIEW v AS ...;
DROP VIEW v;
\end{lstlisting}

	\textbf{Contoh:}
	\begin{lstlisting}[language=SQL]
CREATE VIEW saldo_tinggi AS
SELECT cust_name, balance
FROM customer
WHERE balance > 5000;
\end{lstlisting}

	\textbf{Kegunaan:} sederhanakan query, batasi akses (security), virtual table (tidak simpan data).

	% ==========================================
	\section*{16. Trik Ujian (20 Poin)}

	\textbf{1. NULL.} Jangan tulis \texttt{= NULL}. Selalu \texttt{IS NULL}.

	\textbf{2. N:M.} WAJIB tabel baru. Bukan FK di tabel existing.

	\textbf{3. Cartesian.} \texttt{FROM A, B} tanpa \texttt{WHERE} = bencana. Selalu hubungkan PK-FK.

	\textbf{4. DELETE vs TRUNCATE vs DROP.}
	\begin{itemize}
		\item \texttt{DELETE} — hapus baris, bisa \texttt{WHERE}, bisa rollback
		\item \texttt{TRUNCATE} — semua baris, reset auto-increment
		\item \texttt{DROP} — hapus struktur + data
	\end{itemize}

	\textbf{5. CHAR vs VARCHAR.} \texttt{CHAR} fixed (kode). \texttt{VARCHAR} variabel (nama, email).

	\textbf{6. FK CASCADE.} \texttt{ON DELETE CASCADE} untuk auto-hapus anak.

	\textbf{7. COUNT(*).} \texttt{COUNT(*)} hitung semua baris. \texttt{COUNT(kolom)} hitung non-null saja.

	\textbf{8. GROUP BY.} Semua kolom \texttt{SELECT} non-aggregate wajib ada di \texttt{GROUP BY}.

	\textbf{9. WHERE vs HAVING.} \texttt{WHERE} sebelum grouping. \texttt{HAVING} setelah.

	\textbf{10. LEFT JOIN.} Tetap tampilkan yang tidak cocok (NULL di kanan).

	\textbf{11. LIKE.} Case-insensitive default. Pakai \texttt{BINARY} untuk case-sensitive.

	\textbf{12. TIMESTAMP.} Auto-update dengan \texttt{ON UPDATE CURRENT\_TIMESTAMP}.

	\textbf{13. Compound PK.} \texttt{PRIMARY KEY (a, b)} di akhir \texttt{CREATE TABLE}.

	\textbf{14. Subquery WHERE.} Kalau return banyak nilai, pakai \texttt{IN} bukan \texttt{=}.

	\textbf{15. EXISTS vs IN.} \texttt{EXISTS} lebih cepat untuk korelasi.

	\textbf{16. DECIMAL untuk uang.} \texttt{FLOAT}/\texttt{DOUBLE} tidak presisi. Selisih kecil bisa menumpuk.

	\textbf{17. Aggregate di WHERE.} Tidak boleh. Pakai \texttt{HAVING}. Contoh: \texttt{WHERE COUNT(*) > 5} SALAH.

	\textbf{18. ORDER BY nomor.} \texttt{ORDER BY 1} urut kolom pertama. Portabel tapi tidak jelas.

	\textbf{19. LIMIT.} \texttt{LIMIT 10 OFFSET 20} = ambil 10 baris lewati 20. Setara \texttt{LIMIT 20,10} di MySQL.

	\textbf{20. Alias.} \texttt{SELECT kolom AS nama} wajib \texttt{AS} untuk nama dengan spasi: \texttt{AS 'Nama Lengkap'}.

\end{multicols*}
\end{document}
